\documentclass{kuisthesis}

\renewcommand{\rmdefault}{cmr}  % ローマン体を標準(Computer Modern)に戻す
\renewcommand{\ttdefault}{cmtt}

\usepackage[OT1]{fontenc}
\usepackage[dvipdfmx]{graphicx}
\usepackage{float}
\usepackage{bm}


\jtitle{粘土による直感的操作と生成AIを併用したデザイン探索手法}
\etitle{Design Exploration Method Combining Intuitive Clay Manipulation with Generative AI}
\jauthor{梅木敦成}
\eauthor{Taisei Umeki}
\supervisor{中村裕一 教授}
\date{令和8年2月2日}
\begin{document}
\maketitle
\begin{jabstract}
本研究は、非専門家が持つ曖昧なデザインイメージを具体化するために、手元の粘土の操作による形状入力と生成AIによる解釈・生成を組み合わせた探索的デザイン支援手法を提案する。 
近年、Text-to-Image生成AIの発展により、言語指示による画像生成が容易になった。
しかし、デザインの初期段階において、非専門家が頭の中に描くイメージは不明瞭であることが多く、それを適切にプロンプト（言語）は困難な場合がある。
特に、自然言語は「雰囲気」や「概念」を伝えるには適しているが、「具体的な空間的特徴」や「三次元的な構成」を正確に記述するには限界がある。
さらに、ユーザーが用いる言葉（例：「かっこいい」「やわらかい」）が持つ多義性が、生成AIに正確に目的とする内容を伝える際の問題となる。
加えて、ユーザー自身が自分の用いる言葉の曖昧さに無自覚である場合があり、生成AIの生成結果と想定していたデザインの間に違いがあっても、その原因を特定できずに試行錯誤が長期化するという問題も存在する。

提案手法では、この課題に対し、以下の3段階のプロセスを用いたハイブリッドなアプローチを行う。
(1)「実体入力」：言語化できない形状のニュアンス、ボリューム感、全体のバランスといった非言語的な情報を、可塑性のある粘土模型によって直接入力し、生成AIへ構成情報を伝達する。
また、粘土をこねる身体的行為を通じて、ユーザーの思考を外部化し、アイデア創出を促進する「触覚的思考」を支援する。
(2)「解釈による方向性の先鋭化」：ユーザーの入力した抽象的な形容詞を、LLM（Large Language Model）が具体的な属性（例：金属、有機的曲線、迷彩柄など）へと翻訳して提示しする。
ユーザーが提示された解釈案の中から自身のイメージに近いものを選択することで、言語の曖昧さを解消し、システムとの間でデザインの方向を形成する。
(3)「属性空間での探索」：定義された属性を制約として与えることで、ランダムな生成を防ぎ、ユーザーの意図した構造を保ちながら多様なデザイン候補を探索的に生成・提案する。

被験者を用いた定性的評価の結果、本手法の有用性と課題の双方が明らかになった。
第一に、粘土をこねる身体的行為が、単なる入力手段にとどまらず、ユーザーの思考を先鋭化しアイデア創出を促進する「触覚的思考」を支援することが確認された。
第二に、生成AIによる解釈提案が、ユーザーの潜在的な意図を言語化させる効果を示した。実験では、生成AIが提示した形状定義を選択する過程で、ユーザー自身の使用したが使用語彙の多義性に気づき、言語化する様子が観察された。
一方で、実体入力特有の課題として、粘土のひび割れや凹凸が意図しないテクスチャや模様として過剰に解釈されるノイズの問題や、意図が言語化、明確化した探索の終盤において微細な形状修正を行う際、属性ベースの操作では収束に時間がかかり、ユーザーのフラストレーションにつながる点も確認された。

本研究は、これらの知見に基づき、人間の身体的な直感と生成AIの論理的な表現力を、対話を通じて共存させることで、非専門家の創造性を拡張する新たなデザイン探索プロセスを提示するものである。
\end{jabstract}
\begin{eabstract}
This study proposes an exploratory design support method that inputs clay-based shape to generative AI to help non-expert users materialize vague design ideas. 
Conventional text-to-image generation methods face problems in conveying geometric nuances solely through natural language. 
Furthermore, semantic ambiguity in user prompts often leads to a mismatch between the AI's interpretation and the user's intent, and users are frequently unaware of the gap, leading to undesirable results. 
To address these issues, our method integrates (1) multimodal input using clay photos and text, (2) resolution of semantic ambiguity through "interpretation suggestions" by AI, and (3) exploratory image generation based on the resolved constraints. 
Qualitative evaluation with two participants suggested that hands-on clay manipulation supports ideation and that generative AI suggestions serve as an effective starting point for exploration. 
However, challenges were identified regarding low controllability of generative AI in fine-tuning specific details and the difficulty for users to accurately mold their intended shapes due to lack of modeling skills.
This research presents a novel creative process that harmonizes human physical intuition with generative AI's expressive power through interactive dialogue.
\end{eabstract}
\tableofcontents %目次の出力

% 1章
\section{はじめに}
ものづくりやプロダクトデザインにおける創作活動では、制作の初期段階において「なんとなくこういうものが作りたい」という曖昧なイメージが出発点となることが多い。
しかし専門的な訓練を受けていない非専門家にとって、頭の中のイメージを具体化し、外部へ表現することは容易ではない。
従来はスケッチやCADツールが用いられてきたが、これらは描画能力や操作スキルを要求するため、非専門家にとって難しいプロセスとなっていた。

近年、Text-to-Image を代表とする生成AIの発展により、言語（プロンプト）を入力として多様な画像を生成できるようになった。
生成AIは「描く技術」を担うことが期待されており、アイデアの視覚化を加速させている。
つまり、言語(プロンプト)インタフェースによって、描画や造形のプロセスを生成AIに任せられるようになるのではないかと期待されている。

しかし、生成AIを利用してイメージを具現化するためには、まだ多くの課題が残っている。
本研究ではその課題を以下の二点に整理する。

第一に、言語を用いて「具体的な形状」を伝えることの困難さである。言語は「雰囲気（画風やトーン）」や「概念（物体名や比喩など）」を伝えることには適しているが、「どこをどのように」曲げたいか「どの程度まで」ボリュームを増やしたいかといった空間的な特徴を正確に指示することには困難が伴う。
その結果、ユーザーは言語化できない形状へのこだわりを妥協するか、あるいは偶然生成される出力に依存して試行錯誤を繰り返す状況に陥りやすい。
これは、連続的な『形の空間』と離散的な『言葉の空間』が1対1に対応していない（言葉では表現しきれない形が無数に存在する）ことを意味する。

第二に、言語そのものが持つ「多義性（複数の異なる意味に解釈できる性質）」の問題である。
例えば単に「かっこいい」と指示しても、それが「未来的なかっこよさ」なのか「伝統的なかっこよさ」なのかは人によって異なる。 
重要な点は、ユーザー自身もこの「言葉のあいまいさ」を自覚していない場合が多いことである。 
自分が入力した言葉が多様な解釈を含んでいることに気づかないまま生成を行うため、生成AIが出力する結果に対して「何かが違う」と感じつつも、どう修正すればよいか分からなくなってしまうのである。

こうした「形状記述の困難さ」と「意味解釈の多義性」により、ユーザーが頭の中で描くデザインと、生成AIが生成する出力との間にはどうしてもズレが生じやすく、納得のいく結果を得るまでのプロセスが長期化する傾向にある。
% これは、連続的な『形の空間』と離散的な『言葉の空間』が1対1に対応していない（言葉では表現しきれない形が無数に存在する）ことを意味する。

本研究の目標は、非専門家ユーザーの曖昧で多義的な初期イメージを、システムとの対話を通じて段階的に言語化・具体化し、最終的に明確な要求仕様へと昇華させることである。
そのために、言語化しにくい形状特徴を直接入力できる物理的インタフェースとして粘土模型を用い、その画像とテキストを併用したマルチモーダル入力を採用する。
さらに、LLMによる「解釈の提案」を行うことで、ユーザーに言葉の多義性に気づかせ、曖昧な入力テキストの意味を明確化する枠組みを設計する。
これにより、生成された多様な候補の中からユーザーの意図に沿うものを段階的に絞り込み、納得感のあるデザイン探索を可能にする。

本研究には二つの社会的意義がある。
第一に、高度な造形スキルを持たない人々が自身の内面にあるイメージを正確に外部化できるよう支援することで、個人の創造的な表現活動を促進する点である。
第二に、AIによる自動化が進む現代において、人間が本来持つ「手で考える（触覚的思考）」能力を排除するのではなく、AIによる解釈支援と組み合わせることで、人間の身体的直感とAIの表現力を共存させる新たな創作プロセスを提示する点である。

本論文の構成を述べる。
第2章では、生成AIを用いた物品デザインの背景と、探索的デザイン支援における課題を整理する。
第3章では、提案するデザインプロセスを段階的に説明し、実体入力と画面上の探索を往復する反復ループを示す。
第4章では、ユーザー実験に基づき、提案手法の有用性と課題を定性的に分析する。
第5章では、結論と今後の展望を述べる。



% 2章
\newpage
\section{生成AIを用いた物品のデザイン}
\subsection{デザインにおける生成AIの利用}
% デザインプロセスにおける「探索」の重要性
デザインとは、完成された設計図を最初から出力する作業ではなく、試行錯誤を通じて頭の中の不明瞭なイメージを徐々に明確化していく探索的プロセスである。
特に初期段階では、デザインする人自身も結果のイメージを明確に持っていないことが多い。
そのため、デザインに生成AIを用いる場合には、生成AIからのフィードバックを見ながら意図を更新していく対話的なプロセスが重要となる。

% 言語インタフェースの問題
現在主流のText-to-Image手法は自然言語を入力とし、ユーザーはプロンプトを通じて生成結果を制御する。
言語は「雰囲気」や「概念」といった抽象的情報を伝えることができる一方で、「具体的な形状」や「空間的な構成」といった非言語的情報を正確に指示することは難しい。
例えば、「右側をなめらかに隆起させる」といった形状操作をテキストだけで正確に伝えることは困難である。
この「伝えたい形状イメージ」を「入力手段としての言語」で正確に表現できないというギャップが、意図通りのデザイン生成を妨げる要因となる。

形状以外の属性についてもユーザーが頭の中に描いているイメージを正確にシステムに伝えることは簡単ではない。
つまり、多義的な言葉を用いてしまうと、その意図が正確に伝わらないことがある。
例えば、「かっこいい」車が示唆する具体的な属性は、ユーザーによって大きく異なり、「角ばった」、「車高の低い」、「黒い色調の」など、種々のものが考えられる。
本質的な問題は、ユーザー自身がこの言葉の多義性に無自覚である場合にある。
ユーザーは自身の言葉が意図通りに伝わることを期待して入力を行うため、AIの解釈が自身の期待するデザインとずれた場合、なぜ意図が伝わらないのかを理解できず、適切な修正を行うことが困難となる。

以上の議論から、生成AIをデザインに用いるためには、上記2つの問題を解決するメカニズムが必要である。
そのためには、ユーザーに自身の入力テキストの曖昧さを「意識」させた上でユーザーの意図する特徴を先鋭化し、具体的にシステムへ伝達するためのインタフェースが必要不可欠である。

\subsection{探索的デザイン支援システム}
本研究では、実体をを用いた形状入力と対話的な探索を融合したデザイン支援システムを提案する。
本システムは、粘土による形状入力、LLM(Large Language Model)による多義性解消（解釈提案）、および制約に基づく画像生成探索の三段階で構成される。
これにより、形状と言語の両面からユーザーの頭の中のイメージを具現化させるデザインの実現を目指す。

このために以下のハイブリッドなアプローチをとる。

第一に、物理的インタフェース（粘土）による非言語的な意図伝達を行う。
言語化できない形状のニュアンス（曲がり具合、ボリューム感など）を、粘土という可塑性のある物理メディアを通じて直接入力することで、言語プロンプトの限界を補完する。
第二に、構造的解釈によって多義性を解消する。ユーザーの入力テキストに含まれる曖昧さに対し、生成AIが複数の解釈案を提示する。
ユーザーが自分の意図に合致するものを選択することによってデザインの方向性を決定し、不明瞭なイメージの先鋭化を促す。
第三に、制約に基づく探索的ナビゲーションにより、広大なデザイン空間の中からユーザーの好む領域を探索的に絞り込む。
一回の生成で成果物を出そうとするのではなく、さらに表現を追加したり「もっと鋭く」「もっと有機的に」といった相対的な表現を与えながら探索を継続することを可能にする。

この探索プロセスを実現するために、本研究では検索支援システムFindMeで提案された『制約に基づくナビゲーション』の概念\cite{burke2002}を応用し、生成モデルの探索へと拡張する。
また、意図した範囲内（制約内）で多様な候補を提示することで、ユーザー自身も気づいていなかった好みの発見（発想の拡張）を促すことを狙う。

関連研究として、StyleCLIP\cite{patashnik2021}など、テキスト入力によって画像の特徴を操作する手法が提案されている。
しかし、これらはテキストの意味内容との整合性（意味的な整合性）を重視しており、ユーザー個人の期待するデザインとの整合性や主観的満足度を高める枠組みとは目的が異なる。
本研究は、非専門家が曖昧なイメージを段階的に具体化していくプロセスに焦点を当て、その過程に必要な入力・解釈・探索の支援を統合する点に特徴がある。



% 3章
\newpage
\section{デザインプロセスの設計}
本章では、提案手法を用いたデザインプロセスの設計について具体的に述べる。
本プロセスは、ユーザーの曖昧なイメージを段階的に明確化する「意図の詳細化」と、その意図の範囲内で候補を生成しながら満足できる解を探る「デザイン空間内の探索」を繰り返すことによって進行する。
生成AIは多数の候補を提示できる一方、提示が一貫性のない変形の連続になってしまうと、ユーザーは「どの方向に進んでいるのか」を見失い、試行錯誤のコストが増大する。
そこで本研究では、ユーザーのデザイン意図を先鋭化するプロセスを導入し、探索の方向性を段階的に絞り込みながら詳細化できるように設計する。

% \subsection{デザインプロセスの流れ}


% まず Phase A では、ユーザーは粘土模型の画像と、デザイン意図を表すテキストを入力する。
% 次に Phase B では、入力されたクエリの多義性を解消するために、デザインや接合方法に関する複数の解釈案を提示し、ユーザーが選択することで方向性の合意形成を行う（図\ref{fig:system-flow-abc}）。
% Phase C では、選択された解釈を、画像生成エンジンが処理可能な数値表現（特徴ベクトル）へ変換し、候補画像を生成して提示する。
% Phase D では、提示された画像を見ながらユーザーが特定の属性に対する修正指示（制約）を追加し、その指示を満たす範囲でパラメータ（特徴ベクトル）を更新して再生成を行う（図\ref{fig:system-flow-d}）。
% これにより、ユーザーは制約による『意図の反映』と、AIの生成による『思いがけない発見』を繰り返しながら、意図に近いデザインへと候補を絞り込んでいく。

\subsection{デザインプロセスの流れ}
本システムにおけるデザイン探索は、図\ref{fig:system-flow-entire}に示すように、「入力」「解釈」「変換」「探索」の4つのプロセスを基本単位として進行する。

本手法は、ユーザーの制作段階に応じて2つの開始方法を提供する。
新しく探索を始めたいときは、まだ形が定まっていない段階で抽象的なテキストと粘土画像を入力し、生成AIによる解釈を経てイメージを具体化していく（入力 $\to$ 解釈 $\to$ 変換 $\to$ 探索）。
一方、細部のデザインの調整や探索がしたいときは、ある程度形になった粘土模型を修正した後、その画像を再入力することで、即座にパラメータを推定して探索プロセスへ合流する。
これにより、発想の立ち上げだけでなく、フィードバックループを通じたデザインの洗練も支援する。

入力プロセスでは、ユーザーは粘土模型の画像と、必要に応じてデザイン意図を表すテキストを入力する。 
解釈プロセスでは、入力されたクエリの多義性を解消するために、デザインに関する複数の解釈案を提示し、ユーザーが選択することで方向性の先鋭化を行う（図\ref{fig:system-flow-abc}）。 
変換プロセスでは、選択された解釈を、画像生成エンジンが処理可能な数値表現（特徴ベクトル）へ変換し、候補画像を生成して提示する。 
探索プロセスでは、提示された画像を見ながらユーザーが特定の属性に対する調整指示（制約）を追加し、その指示を満たす範囲でパラメータ（特徴ベクトル）を更新して再生成を行う（図\ref{fig:system-flow-d}）。 
このプロセスでは「特徴ベクトル操作による物体の全体的なデザイン操作」と「一部分のデザインを調整する操作」の2つの操作を提供し、デザインの全体的な印象の調整から局所的なデザイン編集までを一貫した流れで行う。

また、本システムは画面上での完結を前提としない。
画面上で得られた生成画像は、粘土模型を洗練するための参照としても利用される。
ユーザーは生成画像を手がかりに粘土模型を作り直し、再び入力プロセスへ戻ることで、実態と画像表現を往復する探索を行う。

\begin{figure}[H]
  \centering
  \includegraphics[width=0.65\linewidth,keepaspectratio]{/Users/umeking5/university/lab_thesis/flows/flow2_entire.png}
  \caption{本システムの全体フロー}
  \label{fig:system-flow-entire}
\end{figure}

\subsection{デザイン意図の先鋭化}
デザインの初期段階においてユーザーのイメージは曖昧であり、そのままではシステムが処理可能な指示になりにくい。
本節では、入力からデザイン意図の方向を形成するまでのプロセスを述べる。

入力プロセスでは、粘土模型の写真とテキストを扱うが、前述の通りユーザーの目的に応じて処理が分岐する。 
デザイン意図がまだ明確でない場合は、粘土画像と「未来的な椅子」といった抽象的なテキストを入力し、次の言語的な解釈へと進む。 
細部のデザインの調整や探索をする場合では、修正後の粘土画像のみを入力とする。
システムは画像を解析し、現在の粘土模型がどのような属性（丸み、鋭さ等）で構成されているかを逆算して特徴ベクトルを生成し、解釈プロセスをスキップして直接探索プロセス（編集画面）へと遷移する。

% 解釈プロセスでは、非専門家が入力する抽象的なクエリの多義性を解消するために、システムが複数の解釈案を提示する。 
% 例えば「和風にして」という指示クエリに対し、システムは「ヒノキを用いた直線的で簡素な和風」や「漆塗りで装飾的な和風」といった具体的なニュアンスを含んだ解釈案を生成する。
% ユーザーはその中から自身の頭の中のイメージに最も近いものを選択し、デザイン探索の方向性の形成を行う。
% システムは解釈案を一度に3つ提示し、ユーザーはその中から1つを選択するが、いずれも意図に合わない場合は「どれも違う」を選択して再度解釈案を生成することも可能である。
% この時、ユーザーはクエリを自身の語彙の多義性に気づき、クエリを入力することで自身の意図した解釈が再生されるように頑張る。
% ここで「属性」とは「金属的」「角が丸い」等の視覚的特徴であり、属性値はその強度を 0 から 1 の実数値で表したものである。
% また、本研究では多様なデザインを表現するために、独自の「属性空間」を定義した。
% 具体的には、「Material（材質）」「Line（線）」「Surface（表面）」「Color（色彩）」「Pattern（柄）」などの大カテゴリを設け、その下位に「Iron（鉄）」「Organic Curve（有機的曲線）」「Camouflage（迷彩）」といった具体的な属性を配置している。 
% 各属性は0から1の連続値で強度を表し、これらを組み合わせることで数多のデザインバリエーションを表現可能としている。
解釈プロセスでは、非専門家が入力する抽象的なクエリの多義性を解消するために、システムとの対話を行う。
システムは入力されたクエリと粘土画像を分析し、デザインの方向性に関する3つの解釈案を提示する。
例えば「和風にして」というクエリに対し、システムは以下のような具体的なニュアンスを含んだ解釈を生成する。
「解釈A：ヒノキを用いた直線的で簡素な和風」「解釈B：漆塗りで曲線的かつ装飾的な和風」「解釈C：竹を用いた編み込み構造を持つ和風」
ユーザーは提示された選択肢を確認し、自身のイメージに最も近いものを1つ選択する。
もし、提示された3つの解釈がいずれも意図と異なる場合は、「合致しない（No）」を選択する。
この場合、ユーザーは「自分の言葉がシステムにどう誤解されたか」を認識した上で、より適切な言葉を選び直し、クエリや粘土画像を再入力する。
このループを繰り返すことで、ユーザーは自身の言葉の多義性に自覚的になり、意図を正確に伝えるための言語表現を探索する。

本システムにおける「解釈」や「生成」は、独自の「属性空間」に基づいて行われる。
「属性」とは「金属的」「角が丸い」といった視覚的な特徴であり、その強度は0から1の実数値で表される。
本研究では「属性空間」として、「Material（材質）」「Line（線）」「Surface（表面）」「Color（色彩）」「Pattern（柄）」などの大カテゴリを設け、その下位に「Iron（鉄）」「Organic Curve（有機的曲線）」「Camouflage（迷彩）」といった具体的な属性を配置している。
これらを組み合わせることで、数多ののデザインバリエーションを表現可能としている


変換プロセスでは、意図の解釈（自然言語）をシステム内部の数値表現で処理可能な「特徴ベクトル」に変換される。
「特徴ベクトル」とは各属性の値を要素として保持する多次元ベクトルである。 
最終的に、選択された解釈テキストと特徴ベクトルを統合して画像生成プロンプトを構築し、画像生成エンジンへと送られる。


\begin{figure}[H]
  \centering
  \includegraphics[width=0.6\linewidth]{/Users/umeking5/university/lab_thesis/flows/flow2_abc.png}
  \caption{入力から特徴ベクトル生成までのプロセス}
  \label{fig:system-flow-abc}
\end{figure}

\subsection{デザイン意図に基づく探索} 
意図の方向性が定まった後も、細部の調整においては言語化が難しい微細な意図が残る。 
本節では、属性空間内での探索によってデザインを洗練させるプロセスを述べる。

デザインの探索において、各生成結果と特徴ベクトルををセットで「ノード」として定義し、探索木として管理する。
最初に生成した特徴ベクトルを探索木の根ノードとして探索を開始し、生成のたびに子ノードを追加していく。
このとき、ユーザーが現在選択しており、次の操作の基点となるノードを「現在ノード」と呼ぶ。
このように履歴を木構造として保持し、ユーザーが任意の探索段階（ノード）に戻って再探索できるようにすることで、探索が行き詰まった場合のリカバリを容易にする。

全体のデザイン操作をするときは、ユーザーは特徴ベクトルの属性値を増減させて（例：丸っぽさを高める）、物体全体の印象を調整する。
システムはユーザーによる属性値の指定を「等号制約」として受け取り、特徴ベクトルを更新する。 
更新されたベクトル値は、即座に具体的な「形容詞」へと翻訳され、元の解釈テキストに対して追加的な制約として統合される。
システムは初期入力画像と新たなプロンプトを用いて画像生成を行い、候補画像を提示する。
これにより、ユーザーが作成したデザインを大きく崩すことなく、材質や質感、細部の形状といった特定の属性のみを独立して制御する。
結果として、ユーザーは自身の意図する方向へとデザインの探索を進めることができる。

局所的なデザイン操作をするときは、ユーザーは現在ノードの画像上の特定の部位（例：脚）をマスクし、対象となる属性値の増減を指示する。 
システムは入力された変化量に基づき、その属性の強度を段階的に反映したプロンプトを生成する。
生成された新たなプロンプトと現在ノードの画像を用いて画像生成を実行する。
例えば「鋭利さ」を最大まで高めると「より極めて鋭利に（extremely sharper）」という強調表現が適用され、中間値であれば「より鋭利に（sharper）」といった表現が適用される。
これにより、「脚を細くしたいが、どの程度細くするか」といった微細なニュアンスを、対話的に決定することができる。

探索を繰り返す過程では、パラメータ調整を行っても似通った変化が連続してしまい、候補が局所的な解に収束して多様性が失われることがある。
そこで本研究では、ユーザーが明示的に「好ましくない」として棄却したデザインを「クローズドノード」として記録し、探索空間においてこれらのノード周辺を回避する仕組みを導入する。

具体的には、クローズドノードから現在の探索点に対して働く「斥力」メカニズムを実装する。
システムは、カレントノードの特徴ベクトルと、過去のクローズドノード群とのユークリッド距離を計算する。
もしカレントノードがクローズドノードの近傍（一定の閾値以内）に存在する場合、両者の差分ベクトル（クローズドノードからカレントノードへ向かうベクトル）を計算し、これを現在のベクトルに加算する補正を行う。
この処理により、ユーザーが意識せずとも、システムは「過去の失敗作」の類似領域から自動的に離れるようベクトルをシフトさせ、未探索の新しいデザイン空間へと誘導を行う。


\begin{figure}[H]
  \centering
  \includegraphics[width=\linewidth]{/Users/umeking5/university/lab_thesis/flows/flow2_d.png}
  \caption{制約探索によるデザイン洗練プロセス}
  \label{fig:system-flow-d}
\end{figure}

% \subsection{詳細化と探索の繰り返しによるデザイン}
% 以上のプロセスにより、ユーザーはデジタル空間での探索ループを回すことができる。
% 生成画像に対して、全体的な印象や方向性を刷新したい場合は Phase A に戻ってクエリや解釈を再生成し、形状のバリエーションや細部のニュアンスを探索したい場合は Phase D での制約追加を繰り返す。

% さらに、本システムは物理とデジタルの相互作用を含む。
% デジタル上で得られた生成画像が、ユーザーにとって「この方向が良い」「この形は避けたい」といった判断材料となり、粘土模型の修正へとフィードバックされる。
% 例えば椅子のデザインでは、初期の粘土模型で全体の包まれ感を表現し、生成画像を見ながら背もたれの厚みや脚部の広がりを粘土で調整する。
% 修正後の模型を再度入力することで、システムは更新された形状を初期形状として認識し、よりユーザーの意図を正確に反映した提案が可能となる。
% この「物理入力 $\to$ AIによる拡張 $\to$ 物理修正 $\to$ 再入力」というサイクルの反復により、曖昧な初期イメージは実装可能なレベルの具体的デザインへと段階的に昇華される。

\subsection{詳細化と探索の繰り返しによるデザイン}
以上のプロセスにより、ユーザーは探索ループを回すことができる。
生成画像に対して、全体的な印象や方向性を刷新したい場合は入力プロセスに戻ってクエリや解釈を再生成し、デザインのバリエーションや細部のニュアンスを探索したい場合は探索プロセスでの制約追加を繰り返す。

さらに、本システムでは粘土の模型操作と画面上の画像生成を交互に行いながらデザインを進める。
コンピュータ上で得られた生成画像が、ユーザーにとって「この方向が良い」「この形は避けたい」といった判断材料となり、粘土模型の洗練へとフィードバックされる。
例えば椅子のデザインでは、初期の粘土模型で全体の包まれ感を表現し、生成画像を見ながら背もたれの厚みや脚部の広がりを粘土で調整する。
修正後の模型を再度入力することで、システムは更新された形状を初期形状として認識し、よりユーザーの意図を正確に反映した提案が可能となる。
この「実体による入力 $\to$ 生成AIによる拡張 $\to$ 実体修正 $\to$ 再入力」というサイクルの反復により、曖昧な初期イメージは実装可能なレベルの具体的デザインへと段階的に昇華される。



% \subsection{斥力による探索空間の適正化}
% 探索的デザインにおいて、ユーザーが一度「好ましくない」と判断したデザインと類似した候補が繰り返し提示されることは、探索の効率を著しく低下させる要因となる。
% これを防ぐため、本システムでは探索木上の「棄却（Closed）」されたノードから、現在の探索点を自動的に遠ざける「斥力」メカニズムを導入した。

% 具体的には、現在の特徴ベクトル $\bm{v}_{current}$ と、過去に棄却されたノード集合 $C = \{c_1, c_2, \dots, c_n\}$ に含まれる各ベクトル $\bm{v}_{c_k}$ とのユークリッド距離 $d_k = \|\bm{v}_{current} - \bm{v}_{c_k}\|$ を計算する。
% この距離 $d_k$ が一定の閾値 $\epsilon$ を下回る場合、つまり現在の探索点が過去の失敗作に近づきすぎた場合、以下の式に基づいてベクトルを補正する。

% \begin{equation}
%   \bm{v}_{new} = \bm{v}_{current} + \alpha \sum_{k: d_k < \epsilon} w(d_k) (\bm{v}_{current} - \bm{v}_{c_k})
% \end{equation}

% ここで $\alpha$ は斥力の強度係数、$w(d_k)$ は距離に応じた重み関数（距離が近いほど強く反発する）である。
% 上式の第2項は、棄却されたベクトル $\bm{v}_{c_k}$ から現在のベクトル $\bm{v}_{current}$ へ向かう方向、すなわち「失敗から離れる方向」へのベクトル加算を意味する。
% この処理により、過去の不適切な解の周辺を再訪することを防ぎ、ユーザーが意識せずとも常に新しいデザイン空間へと探索が誘導される効果が得られる。





% 4章
\newpage 
\section{実験例と考察} 
本章では、提案手法（粘土による形状入力と生成AIを活用した探索）が、デザインの専門知識を持たないユーザーに満足できるデザインを得るための有効な支援となっているかを検証するための実験について述べる。

\subsection{システム構成} 
本システムは、Webブラウザ上で動作するアプリケーションとして実装した。 
ユーザーインターフェースの実装にはPythonのWebフレームワークであるStreamlitを採用している。 
バックエンド処理にはOpenAI社のAPIを使用しており、ユーザーの言語指示や入力画像の解釈には「GPT-4o」、画像の生成には「GPT-Image-1」、および画像編集機能（Inpainting）を活用している。 
これらを統合することで、スライダー操作やマスク描画といったユーザーの直感的な操作を、リアルタイムに生成AIへの生成指示（プロンプト）へと変換する仕組みを構築した。


\subsection{実験方法} 
提案手法の有効性を検証するため、以下の3つの条件を設定した。

\begin{enumerate} 
  \item \textbf{提案手法}: 粘土画像とテキストを入力後、解釈案の選択と属性制約による探索を行う手法。 
  \item \textbf{比較手法A（粘土＋テキストのみ）}: 粘土画像とテキストで初期生成を行った後、テキストプロンプトのみで画像の反復生成（Image-to-Image）を行う手法。 
  \item \textbf{比較手法B（テキストのみ）}: テキストのみを入力し、反復生成時も直前の生成画像を入力に含めず、テキストのみで生成（Text-to-Image）を行う手法。 
\end{enumerate}

実験の評価は、主に作業中の発話および実験後の半構造化インタビューに基づいた定性的な分析により行う。 
被験者は3名（以下、被験者P、被験者Q、被験者R）とし、日用品を対象とした2つの課題（「椅子をデザインしてください」「カップをデザインしてください」）に取り組んでもらった。 
課題には正解画像を設定せず、被験者が満足するまで探索を繰り返してよいものとし、時間制限は設けなかった。 
なお、本システムは生成されたデザインを手がかりにユーザーが実体(粘土)を修正するプロセスを前提としている。 
そのため、物理的な即時修正が困難な「色彩」や「材質」といった属性は除外し、フィードバックループの検証に適した「形状（Shape）」に関する53属性(表\ref{tab:attribute-compact})に焦点を絞って実施した。 
また、被験者Pの実験については部分編集機能の実装前であったため、全体編集のみを評価対象としている。

\begin{table}[H]
  \centering
  \small
  \caption{定義した属性空間の構成（抜粋または全一覧）}
  \label{tab:attribute-compact}
  \begin{tabular}{p{3cm} l l}
    \hline 
    \textbf{Category} & \textbf{Attribute} & \textbf{Description} \\
    \hline 
    \textbf{1. Form} & cubic & 立方体・箱型 \\
     & cylindrical & 円柱状 \\
     & (省略...) & ... \\
    \hline 
    \textbf{2. Line} & straight & 直線的 \\
     & organic\_curve & 有機的曲線（不規則なうねり） \\ 
     & (省略...) & ... \\
    \hline 
    \textbf{3. Balance} & asymmetrical & 非対称 \\
     & thick & 肉厚 \\
     & (省略...) & ... \\
    \hline 
    \textbf{4. Edge} & sharp & 鋭い \\
     & round & 丸みを帯びた \\
     & (省略...) & ... \\
    \hline 
  \end{tabular}
\end{table}

実験条件の詳細と設定意図を表\ref{tab:conditions}に示す。

\textbf{提案手法}は、粘土の形状操作と、属性操作による詳細な制御の双方を用いる条件である。

\textbf{比較手法A}は、提案手法から「属性制約（特徴ベクトル）」を除外した構成に相当する。 
直前の生成画像を次の入力として用いることで形状の維持は図られるものの、言語指示のみによる微調整がいかに全体構造の崩れを招きやすく、意図した修正が困難であるか（制御性の低さ）を検証するために設定した。

\textbf{比較手法B}は、提案手法から「粘土による実体入力」を除外した構成に相当する。 
実体的な手がかりがない場合、多義的な言語指示のみではユーザーの意図する具体的な形状特徴を生成結果に反映させることが困難であることを検証するために設定した。

\begin{table}[H]
  \centering
  \caption{実験条件の概要}
  \label{tab:conditions}
  \begin{tabular}{|p{0.18\linewidth}|p{0.36\linewidth}|p{0.36\linewidth}|}
    \hline
    条件 & 初期入力 & 反復時の入力・操作 \\
    \hline
    提案手法 & 粘土画像＋テキスト & 解釈案の選択、特徴ベクトル生成、属性制約による探索（必要に応じて部分編集／合成） \\
    \hline
    比較A & 粘土画像＋テキスト & テキストプロンプトによる反復生成（直前の生成画像を次の入力に含める） \\
    \hline
    比較B & テキストのみ & テキストプロンプトによる反復生成（直前の生成画像は入力しない） \\
    \hline
  \end{tabular}
\end{table}

実験後インタビューは半構造化形式で行った。
本研究では「非専門家はイメージが曖昧であり、言語化が困難である」という課題に対し、システムが「解釈の提案」を通じてその具体化を支援すると仮定している。
そのため、既存の操作感に関する質問に加え、ユーザーの頭の中のイメージがプロセスのどの段階で明確になったか、生成AIの提示によって潜在的な意図が言語化されたかといった「認知的な変化」に焦点を当てた質問項目を追加した。

具体的には、手を動かすことの影響（Q1）、初期出力とのギャップ（Q2）に加え、生成AIの解釈による言語化の気づき（Q3）、自身のイメージの曖昧さへの自覚（Q4）、システムからの提案がデザインに与えた影響（Q5）、修正時の制御感（Q6）、言語化の困難さ（Q7）、属性選択方式への評価（Q8）、そして最終的なイメージの鮮明化と納得感（Q9, Q10）について聴取した。
インタビュー項目の構成とねらいを表\ref{tab:interview}に示す。

\begin{table}[H]
  \centering
  \caption{実験後インタビュー項目の概要}
  \label{tab:interview}
  \begin{tabular}{|p{0.08\linewidth}|p{0.25\linewidth}|p{0.61\linewidth}|}
    \hline
    ID & 項目 & ねらい \\
    \hline
    Q1 & 身体性の影響 & 手を動かす（粘土造形）ことが、発想やイメージの変化に与える影響を把握する。 \\
    \hline
    Q2 & 初期ギャップ & 最初の生成結果に対する「意図が伝わっていない」感覚が、どの条件で強いかを把握する。 \\
    \hline
    Q3 & 言語化による発見 & 生成AIが提示した「解釈（言葉）」を見ることで、自身の作りたかったものが言語化されたか、あるいは新たな発見があったかを確認する（多義性の解消）。 \\
    \hline
    Q4 & 曖昧さの自覚 & 多様な解釈案や予想外の生成画像を通じて、自身の中で「まだ決まっていなかった部分（未決定領域）」に気づけたかを確認する。 \\
    \hline
    Q5 & 提案の影響 & システムが提示した属性や解釈が、最終的なデザインの形状変更やアイデア拡張にどう寄与したかを把握する。 \\
    \hline
    Q6 & 微調整の制御感 & 「ここをこう直したい」という修正場面において、属性制約による操作とテキスト指示による操作のどちらが意図通りに制御できたかを比較する。 \\
    \hline
    Q7 & 言語化の壁 & 比較条件（テキスト指示）において、「イメージはあるのに言葉が見つからず妥協した」場面の有無と内容を把握する。 \\
    \hline
    Q8 & システムの制約 & 「属性リストから選ぶ」という形式が制約（不自由さ）として感じられた場面や、自由記述への要望を把握する。 \\
    \hline
    Q9 & イメージの鮮明化 & 曖昧だったイメージが、「粘土操作時」と「生成AIによる視覚化・言語化時」のどの段階で確信（正解像）へと変わったかを特定する。 \\
    \hline
    Q10 & 最終的な納得感 & 最終成果物に対し、「自分だけのこだわり」を反映できたと感じる条件とその理由（形状のニュアンス、プロセスの納得感）を把握する。 \\
    \hline
  \end{tabular}
\end{table}

定性的評価として、これらのインタビュー回答と作業中の発話を内容分析し、操作感、制御可能性、満足度、および発想支援効果（イメージの具体化プロセス）の観点で整理した。

\subsection{結果の説明}
本節では、実験後インタビューおよび実験中の発言記録をもとに、提案手法の有用性と課題を整理する。
結果は、（1）実体的形状入力の影響、（2）イメージの具体化と言語化、（3）探索的発見と制御性の課題、の三つの観点からまとめる。

\subsubsection{実体の形状入力（粘土）の影響}
粘土を用いた実体の形状入力は、発想の誘発と構成の伝達という観点で肯定的に捉えられた。
被験者Qは、テキストのみでの試行錯誤では「いや違うんやけどな」というストレスが生じやすい一方、粘土では手を動かしながら考えられるため負荷が低いと述べている。
また「新しいものができるのはやっぱりこねてるとき」という発言は、粘土操作が単なる入力手段ではなく、アイデア創出を支える活動そのものになっていることを示唆する。

一方で、粘土による造形はユーザーのスキルに依存する側面がある。
被験者Pは「頭でイメージしたことを手作業で表現することが難しい」と述べたが、同時に言葉で説明するよりも「手で作って伝えたほうが構成要素を伝えるのが楽」とも述べており、全体構成を生成AIに伝える手段としては有効であると評価された。
% 被験者Rは粘土を入力しても生成AIに構成が反映されない
ただし、被験者Qの実験で見られたように、粘土表面の意図しない割れ目や指の跡による陰影が画像特徴として強く反映され、生成結果にノイズとして現れる課題も確認された。
図\ref{fig:mug-example}はマグカップタスクにおける生成例であるが、粘土表面の微細な凹凸や意図しない「シワ」が、画像生成エンジンによって入力クエリの「溶けたような複雑なテクスチャ」や、「装飾的な文様」として過剰に解釈されている。
被験者Pが指摘した「不要なディテール」や「不気味さ」は、このような実体入力のノイズが直接的に生成結果へ反映されたことに起因すると考えられる。

\begin{figure}[H]
  \centering
  \includegraphics[width=0.46\linewidth]{/Users/umeking5/university/lab_research/truecoding/data/images/initial_20260126_190138_a11ae3d9.jpg}
  \hfill
  \includegraphics[width=0.46\linewidth]{/Users/umeking5/university/lab_research/truecoding/data/images/generated_20260126_190711_4562329e.png}
  \caption{実体のノイズが生成AIの解釈に影響を与えた例}
  \label{fig:mug-example}
\end{figure}

\subsubsection{イメージの具体化と言語化}
本研究の主眼である「曖昧なイメージの具体化」について、被験者Rの発言から顕著な効果が確認された。

第一に、生成プロセスを通じた「正解像の発見」である。
被験者Rは当初「おしゃれなカップ」という非常に曖昧なイメージで開始した。システムはこれに対し「上部を広げて円錐状にし、全体的に曲線的なエッジを持たせて、優雅さを表現する」という解釈案を提示し、ユーザーはこれを選択した。
しかし、実際に生成された「ずんぐりむっくりなカップ」を見た瞬間、ユーザーは「こういうことじゃない」という強い違和感を抱いた。
この違和感がきっかけとなり、ユーザーは自身が求めていたものが一般的なマグカップではなく「ティーカップ」であることを発見した。
これは、生成AIによる視覚的な提示が、ユーザー自身の潜在的な選好を鏡のように映し出し、曖昧だったゴールを明確化させる「反射的実践」として機能したことを示している。
第二に、「未決定領域（曖昧さ）」の自覚である。
被験者Rは椅子の生成において「背もたれのデザインが理想通りにいかない」と悩みつつも、インタビューでは「生成前は背もたれの形までは考えていなかった（自分でも曖昧だった）」と振り返っている。
生成AIが多様な解釈案を提示することで、ユーザーは自身が「決めたつもりで決めていなかった部分」を自覚し、そこから新たな検討（言語化）を始めるきっかけを得ていた。

第三に、言語化の支援である。
被験者Rは「形状を網羅する属性セットを作るのは難しい」と指摘しつつも、「提示された解釈案の中に『これこれ！』となる良い言葉があった」と述べている。
これは、システムによる解釈提案が、ユーザーの語彙不足を補い、自身の意図を他者（生成AI）に伝えるための言語獲得を支援したことを示唆する。

\subsubsection{探索的発見と制御性の課題}
属性制約による探索では、思いがけない発見の効果が観察された。
被験者Qは「有機的曲線」という属性を探索中にデザインに追加した際、「頭になかった視点を示してくれた」「ぐるぐるにしてくれたのが面白かった」と述べ、自身の想定外の候補を好意的に受け止めた。
この点は、本研究が目指す「探索的デザイン支援」の有用性を支持する。

しかし、明確な意図を持った修正（収束フェーズ）においては課題も浮き彫りとなった。
被験者Pは「場所の指定ができない」と述べ、全体に対する属性操作では局所的な意図を伝えにくいことを指摘した。
また、被験者Qも「明確に言語化された意図があってコントロールしたいなら制約による探索はやってられない」「自分の言葉で生成した方が早い」と述べ、デザインを絞り込む局面においては操作が冗長に感じられると評価した。

さらに、言語的な解釈と生成結果のズレも確認された。
被験者Pの「自分の想像する『くびれ』と生成AIが対象とする『くびれ』が違う」という発言や、被験者Qの「長い文章でも結局出力は文章から想像される形状イメージとはズレる」という指摘は、LLMを介した間接的な操作の限界を示している。
今回は形状に関する属性を用いたが、全ての形状特徴を網羅するための属性セットを構築することは困難であり、ユーザーの意図と生成結果との間にギャップが生じやすい。
これを受けて被験者Qは「部分的な修正に関しては、解釈を介さずに制約（パラメータ）で直接操作したい」という意向を示しており、探索フェーズ（発散）と修正フェーズ（収束）で適切なインタラクション手法を使い分ける必要性が示唆された。

\subsubsection{比較手法との差異}
比較A（粘土＋テキストのみ）は、明確な修正意図がある場合に制御しやすい傾向が見られた。
被験者Pは「比較手法Aの方が良かった」と述べ、特に場所の指定など局所的な指示を言葉で与えられる点を評価した。

比較B（テキストのみ）は、デザインの不確定要素が多く、発想を広げるという観点で有用であった。
被験者Pは「ポンと出してくれるからたまにいい」と述べており、例えば「日本刀」のような粘土での造形が難しい具体的なモチーフを、造形の手間をかけずに短時間で出現させられる点は、探索の初期段階において有用であることが示唆された。
ただし、形状の絞り込みには課題が残り、結果として比較Aや提案手法へ移行して調整する方が良いという意見も得られた。
また、被験者からは「適切な言葉を探すのが疲れる」「諦めて妥協した」という発言が相次いだ。
視覚的なコンテキスト（粘土）がない状態でのプロンプト探索は、ユーザーに適切な言語表現を探すための思考やストレスを強いることが確認された。

\subsection{考察}

実験結果、特に被験者Rの事例からは、提案手法における「解釈プロセス」が、単なるパラメータ設定の手間ではなく、ユーザーの曖昧なイメージを段階的に具体化するための重要な手がかりとして機能していることが示唆された。

デザインの初期段階において、非専門家のユーザーは「おしゃれ」「かっこいい」といった抽象的で多義的な形容詞を用いて検索を行う傾向がある。
しかし、これらの言葉は具体的なデザイン情報を含まないため、従来のText-to-Image（比較手法B）ではAIの解釈に依存したランダムな出力になりやすく、ユーザーの意図したデザインへの収束が困難であった。

これに対し、本システムはこれらの形容詞を具体的な属性、本実験では具体的な「幾何学的特徴」へと翻訳して提示する機能を持つ。
実際に被験者Rのマグカップ課題のログを確認すると、ユーザーの初期クエリである「おしゃれ」に対し、システムは「[2] 直線的で鋭角」「[3] 円錐状で曲線的」といった形状的な定義を与えている。
ユーザーはこの選択肢の中から[3]を選択することで、「自分にとっての『おしゃれ』とは、優雅で曲線的なものである」という定義を、システムとの対話を通じて事後的に具体化させたと考えられる。
このプロセスは、ユーザーが脳内で漠然と感じている「雰囲気」と、それを実現するための「実体の形状」との間のギャップを、生成AIが提示する選択肢によって埋める行為と言える。

さらに興味深い点は、システムとの対話が、当初のクエリそのものの修正（概念の再定義）を促した点である。

被験者Rは当初「マグカップ」というコンセプトで探索を開始したが、システムが提示した解釈案（曲線的で優雅な形状）を選択し、生成された画像を確認する過程で、自身の求めていたものが実用的な「マグカップ」ではなく、より装飾的な「ティーカップ」であることに気づいた。
その結果、次のターンではクエリを「英国風」へと変更し、コンセプトを「ティーカップ」へと修正している。

この際、システムは「英国風」という文化的文脈に対し、以下のような具体的な形状解釈を提示している。
\begin{itemize}
  \item 解釈案 [1]: 直線的かつ対称、ハンドルを曲線的にする。
  \item 解釈案 [2]: リブを加えてエンボス模様をつけ、装飾的にする。
  \item 解釈案 [3]: 縁を優雅に波打たせ、ハンドルを細く繊細にする。
\end{itemize}

ユーザーはここで[3]を選択しており、最終的に「英国風＝縁が波打った繊細な形状」という具体的なデザイン要件を決定した。
比較手法B（テキストのみ）においても試行錯誤による修正は可能だが、それはAIが出力した結果を見て良し悪しを判断する受動的なプロセスになりがちである。その場合、ユーザーは「英国風」という言葉が内包する形状の多義性に無自覚なまま、偶然出力された形状を受け入れることになる。
一方、提案手法では、「英国風」という言葉に対し複数の形状解釈が並列に提示されるため、ユーザーは「自分の言葉がいかに多様な意味を含んでいるか」を自覚させられる。
その上で、提示された選択肢から一つを能動的に選ぶというプロセスを経ることで、ユーザーは無自覚な多義性を解消し、意図した形状特徴をデザインに確実に反映させることができたと考えられる。

以上のことから、本システムにおける「粘土入力 $\to$ 解釈選択 $\to$ 生成」というループは、単なる画像の修正作業ではなく、ユーザー自身の内部にある「曖昧なイメージの解像度を高め、言語化する認知プロセス」を支援していると結論付けられる。




% 5章
\newpage
\section{結論と展望}

本研究では、非専門家が抱く曖昧なデザインイメージを具体化するために、粘土による形状操作と、生成AIによる解釈・生成を組み合わせた探索的デザイン支援手法を提案した。
本手法は、言語だけでは伝えきれない「形状の空間的特徴」を粘土で入力し、言語だけでは自覚しにくい「言葉の多義性」を生成AIによる解釈提案（属性定義）で補うことで、ユーザーの意図を汲み取りながら段階的にデザインを具体化することを目指した。

ユーザー実験を通じた定性的評価の結果、以下の知見が得られた。

第一に、実体を用いた入力手段（粘土）の有用性である。
粘土をこねるという身体的な行為は、単なる形状入力の手段にとどまらず、言語化に行き詰まったユーザーの思考を助け、アイデアの創出を促進する効果が確認された。
これは、コンピュータツールのみでは抜け落ちがちな「手で考える」プロセスを生成AIに取り入れたことの意義を示している。

第二に、生成AIによる「解釈の提案」がもたらす認知的な支援効果である。
被験者Rの事例が示したように、生成AIが曖昧な言葉（例：「おしゃれ」）を具体的な幾何学的特徴へと翻訳して提示することで、ユーザーは自身の言葉の多義性に気づき、意図の言語化の手助けをすることができた。


第三に、現状の課題である。
実体入力には、ユーザーの造形スキルへの依存や、粘土のひび割れ等がノイズとして解釈される問題が伴うことが明らかになった。
また、形状に関しては網羅的な属性セットを用意するのは難しく、ユーザーの意図と生成結果との間にズレが生じやすいこと、解釈から想起される生成結果がユーザーの期待と異なる場合があることも確認された。


本研究の更なる発展に向け、以下の展望が挙げられる。


本実験では、粘土模型との整合性を考慮し、検証対象を「形状（Shape）」に関する属性に限定した。
しかし、第3章で定義した「属性空間」は、本来「材質（Material）」「色彩（Color）」「表面の質感（Texture）」「模様（Pattern）」といった視覚的属性を包括する概念である。
「形状」においては実体入力（粘土）が強力な手がかりとなるが、一方で「材質」や「表面」といったカテゴリの属性こそ、言語による多義性が生じやすい領域である（例：「温かみのある」という言葉が、暖色を指すのか、木材の質感を指すのかは曖昧である）。
したがって、今後は形状以外の属性も実験条件に組み込み、本手法の提案する「解釈と合意形成」のプロセスが、色彩や材質の決定においてもユーザーのデザイン意図決定を支援できるかを検証する必要がある。


粘土特有のノイズ（シワや割れ目）が意図しないテクスチャとして解釈される問題を解決するために、画像処理によるスムージングや、深度センサを用いた形状情報の抽出精度の向上が求められる。
また、粘土だけでなく、積み木や既存の日用品を「見立て」て入力するなど、ユーザーのスキルに依存しない多様な実体入力手段への対応も検討すべきである。

現状のシステムは2D画像を出力するため、物理的に成立しない矛盾した形状が生成される場合や、側面の詳細が確認できない場合がある。 
粘土という3次元入力を活かすため、生成結果を即座に3Dモデルとして提示する方向性が考えられる。 
これにより、ユーザーは全角度からデザインを確認できるだけでなく、生成されたモデルを3Dプリントして物理的な芯材として利用するなど、コンピュータ上での検討と現実での制作を滑らかに行き来する新しい制作フローが実現できる。


本研究が提案した「粘土の操作」と「生成AIの論理的な解釈」を融合させるアプローチは、専門的なスキルを持たない人々が、自身の内面にあるイメージを形にし、創造的な活動を楽しむための新たな一歩になると期待する。

\acknowledgments
本研究を遂行するにあたり、終始多大なるご指導とご鞭撻を賜りました指導教員の中村裕一先生に心より感謝申し上げます。

また、研究の進捗において貴重なご助言とフィードバックをいただいた近藤一晃先生、下西慶先生に深く感謝いたします。

最後に、本研究の実験において多大なるお時間をいただき、熱心にご協力くださった実験協力者の皆様、並びに日頃より有益な議論を交わし研究生活を支えてくださった研究室の皆様に厚く御礼申し上げます。
%
\bibliographystyle{kuisunsrt} % または kuissort
\bibliography{kuisthesis_jp}
%
% 付録があれば \appendix/\Appendix に続いてここに記す。
\end{document}

